\begin{abstract}
Reliable model validation in chemical data analysis depends on how observations are partitioned into training and test sets. In molecular property prediction, random splitting can retain chemically close train/test neighbours, whereas scaffold splitting removes framework overlap but may introduce target-distribution shift. We present a reproducible chemometric framework for diagnosing these effects across six public classification and regression datasets. Logistic or ridge regression, random forest, and XGBoost models were evaluated from Morgan fingerprints under random, ordinary scaffold, and target-balanced scaffold protocols over 20 fixed seeds. The target-balanced procedure was formalized as a stochastic greedy search that minimizes normalized test-size deviation and target-mean deviation while preserving scaffold disjointness. Its target balance was compared with 5,000 random scaffold assignments per dataset. Predictive effects were quantified through paired generalization-gap reductions, bootstrap 95\% confidence intervals, Wilcoxon signed-rank tests, and Holm correction. Target balancing reduced absolute target shift in five of six datasets and, in those datasets, outperformed approximately 91.7--97.2\% of random scaffold assignments on target balance. Its predictive effect was nevertheless heterogeneous: among 18 dataset--model comparisons, Holm-adjusted tests identified nine smaller gaps, six larger gaps, and three inconclusive changes. Chemical similarity, scaffold concentration, and model rankings also varied with the partition. These results show that controlling one split-induced confounder does not create a universally easier or more valid benchmark. Molecular benchmarks should therefore report partition diagnostics together with predictive scores.
\end{abstract}

\textbf{Keywords:} chemometrics; model validation; molecular property prediction; scaffold splitting; benchmark diagnostics.